\section{Introduction}

Ce laboratoire a pour objectif la mise en œuvre pratique du bus de terrain CAN
entre deux cartes Arduino Mega 2560 équipées d'un Shield CAN (MCP2515 +
MCP2551) et d'une carte d'entrées/sorties HEIG-VD. Le but est d'écrire un
protocole applicatif permettant de recopier sur chaque carte les entrées
(tout-ou-rien et analogiques) lues par la carte jumelle.

Le travail se décompose en trois parties :
\begin{itemize}
    \item \textbf{Programmation} du contrôleur CAN MCP2515 à l'aide de la
    bibliothèque \texttt{BusTerUtils} : configuration des masques/filtres de
    réception, gestion des messages analogiques et tout-ou-rien, lecture des
    DIP-switches pour l'adressage.
    \item \textbf{Mesures sur un bus libre} à \num{1}~Mbps : analyse d'une
    trame, comportement lorsque le câble est débranché, mesure de la latence
    et de la gigue entre deux cartes, puis transmission d'un signal
    analogique et sinusoïdal.
    \item \textbf{Mesures sur un bus chargé} à \num{200}~kbps avec deux paires
    de cartes partageant le même bus : étude de l'impact du trafic sur la
    latence, la gigue et la qualité de reproduction d'un signal sinusoïdal.
\end{itemize}

Les mesures sont effectuées à l'oscilloscope Tektronix et les oscillogrammes
sont intégrés au rapport.
